% 图 2.5 经典顺磁体磁化：朗之万函数 L(y)=coth(y)-1/y；虚线为渐近线与原点处切线 y/3
\documentclass[tikz,border=5pt]{standalone}
\usetikzlibrary{arrows.meta}
\usepackage{amsmath}
\usepackage{bm}
\usepackage{fontspec}
\setmainfont{Times New Roman}
\usepackage{pgfplots}
\pgfplotsset{compat=1.18}
\begin{document}
\begin{tikzpicture}[line width=0.8pt,>=Stealth]
\begin{axis}[
  width=10.5cm,height=6.4cm,scale only axis,
  axis lines=middle,
  axis line style={-Stealth,line width=0.9pt},
  xmin=-4.8,xmax=5.4,ymin=-1.30,ymax=1.30,
  xtick={-4,-3,-2,-1,0,1,2,3,4},
  ytick={1,-1},
  tick style={line width=0.7pt},
  xlabel={$\mu B/k_{\mathrm{B}}T$},
  x label style={at={(axis cs:5.4,0)},anchor=west},
  ylabel={$M/M_{\mathrm{s}}$},
  y label style={at={(axis cs:-0.25,1.30)},anchor=south},
  clip=true]
  % 渐近线
  \addplot[dotted,line width=0.7pt,domain=-4.8:5.4] {1};
  \addplot[dotted,line width=0.7pt,domain=-4.8:5.4] {-1};
  % 原点处切线 L ~ y/3
  \addplot[dotted,line width=0.7pt,domain=-4.4:4.4] {x/3};
  % 朗之万函数（分两段避开 y=0 数值奇点）
  \addplot[very thick,line width=1.15pt,domain=-4.8:-0.02,samples=181]
    {(exp(2*x)+1)/(exp(2*x)-1) - 1/x};
  \addplot[very thick,line width=1.15pt,domain=0.02:4.8,samples=181]
    {(exp(2*x)+1)/(exp(2*x)-1) - 1/x};
\end{axis}
\end{tikzpicture}
\end{document}
